\documentclass[11pt]{article}
\usepackage[margin=1in]{geometry}
\usepackage{parskip}
\usepackage{amssymb,amsmath,amsfonts,verbatim,amsthm, amscd}
\usepackage[dvipsnames]{xcolor}
\usepackage[breakable,skins]{tcolorbox}
\usepackage{graphicx,tikz-cd,adjustbox}
\usepackage{float} % breqn can sometimes conflict with amsmath, be mindful
\usepackage{scalerel}
\usepackage{stackengine,wasysym}
\usepackage{lipsum}
\usepackage[backend=biber]{biblatex} % Using biber backend is recommended

\newtcolorbox{mybox}[2][]{
    arc=0mm, enhanced, frame hidden, breakable
}
\usepackage{mathtools}

\usepackage{subcaption}% <-- added

\usepackage{array} % core package
\usepackage{newtxtext,newtxmath} %fuente

\usepackage[cal=boondoxo,scr=euler]{mathalfa}
\usepackage[unicode=true]{hyperref} % permits navigating on the pdf
\hypersetup{pdfencoding=auto}
\usepackage{cleveref} % smart referencing
\usepackage{caption} %

% Define math commands
\newcommand{\calL}{\mathcal{L}}
\newcommand{\calH}{\mathcal{H}}
\newcommand{\bz}{\mathbb{Z}}
\newcommand{\lt}{\normalfont\text{lt}}
\newcommand{\calA}{\mathcal{A}}
\newcommand{\ind}{\text{\normalfont Ind}}
\newcommand{\stab}{\text{\normalfont Stab}}
\newcommand{\hilbM}{\normalfont\text{\underline{H}}_\text{M}}
\newcommand{\aughilbM}{\normalfont\text{H}_\text{M}}
\newcommand{\rank}{\normalfont\text{rk}}
\newcommand{\flats}{\normalfont\text{Flats}}
\newcommand{\flags}{\normalfont\text{Flags}}
\newcommand{\matM}{\normalfont\text{M}}
\newcommand{\hilbMmodF}{\normalfont\text{\underline{H}}_\text{M/F}}
\newcommand{\M}{\normalfont\text{M}}
\newcommand{\N}{\normalfont\text{N}}
\newcommand{\e}{\normalfont\text{e}}
\newcommand{\rk}{\normalfont\text{rk}}
\newcommand{\corank}{\normalfont\text{crk}}
\newcommand{\crk}{\normalfont\text{crk}}

\newcommand{\qbinom}[2]{\begin{bmatrix} #1 \\ #2 \end{bmatrix}_q}

% Theorem environments
\newtheorem{theorem}{Theorem}
\newtheorem{corollary}[theorem]{Corollary}
\newtheorem{lemma}[theorem]{Lemma}
\newtheorem{proposition}[theorem]{Proposition}
\newtheorem{conjecture}[theorem]{Conjecture}
\newtheorem{example}[theorem]{Example}

\theoremstyle{remark}
\newtheorem*{remark}{Remark} % Unnumbered remark
% \theoremstyle{definition} % For proof environment styling - amsthm handles this with \begin{proof}
% \newtheorem*{proof_env}{Proof}

\newcommand{\luis}[1]{{\color{blue} \sf $\clubsuit$ Luis: [#1]}}
\newcommand{\nutan}[1]{{\color{violet} \sf $\clubsuit$ Nutan: [#1]}}
\newcommand{\jacob}[1]{{\color{Bittersweet} \sf $\clubsuit$ Jacob: [#1]}}

\title{\texorpdfstring{Notes (Inverse Kazhdan-Lusztig Polynomials under Deletion)}{Notes (Inverse Kazhdan-Lusztig Polynomials under Deletion)}}
%\author{Nutan Nepal}
%\date{\today}

\begin{document}
\textbf{Notes: Deletion Formula for $Y_M(x)$}\hfill {\small{\today}}

\hrulefill % Horizontal line after title block

Let $M$ be a matroid and $e \in M$. The polynomial $\hat{Y}_M(x)$ is defined as
\[
    \hat{Y}_M(x)
    = \sum_{F \in \calL(M)} \left( x^{\corank(F)}\mu[F,E] \right) \cdot \hat{Q}_{M^F}(x).
\]
For $F\in\calL(M)$, we have,
\[
    \hat{Y}_{M^F}(x)
    = \sum_{G \leq F} \left( x^{\rank(F)-\rank(G)}\mu[G,F] \right) \cdot \hat{Q}_{M^G}(x)
\]
and
\[
    \hat{Y}_{M/G}(x)
    = \sum_{F \geq G} \left( x^{\rank(E)-\rank(F)}\mu[F,E] \right) \cdot \hat{Q}_{M^F_G}(x).
\]


We want to find a formula for $Y_M(x):=(-1)^{\rank(M)} \hat{Y}_M(x)$ in terms of $Y_{M\setminus e}(x)$.

\section*{Suggestion 1}
Let $\mu_F := \sum_{G \leq F} x^{\rank(G)}\mu[G,F]\cdot G$. Then, writing $\mu_M$ in Kazhdan-Lusztig basis, we have
\begin{align*}
    \mu_M & = \sum_{F \in \calL(M)} \left(\left( x^{\rank(F)}\mu[F,E] \right) \cdot \sum_{G \leq F} x^{\rank(F)-\rank(G)}\hat{Q}_{M_G^F}(x^{-2})\cdot \zeta^G \right) \\
          & = \sum_{F \in \calL(M)} \sum_{G \leq F} x^{2\rank(F)-\rank(G)}\mu[F, E]\cdot \hat{Q}_{M_G^F}(x^{-2})\cdot \zeta^G                                         \\
          & = \sum_{G \in \calL(M)} \sum_{F \geq G} x^{2\rank(F)-\rank(G)}\mu[F, E]\cdot \hat{Q}_{M_G^F}(x^{-2})\cdot \zeta^G                                         \\
          & = \sum_{G \in \calL(M)} x^{-\rank(G)}\left(\sum_{F \geq G} x^{2\rank(F)}\mu[F, E]\cdot \hat{Q}_{M_G^F}(x^{-2})\right)\cdot \zeta^G                        \\
\end{align*}

The exponent of $x$ in the above expression should be $-2\corank(F) = 2\rank(F) -2\rank(M)$ for us to get the
expression for $\mu_M$ in terms of $\hat{Y}_{M/G}(x)$.

Redefining $\mu_F$ as $\sum_{G \leq F} x^{\rank(G) - 2\rank(F)}\mu[G,F]\cdot G$, we have
\begin{align*}
    \mu_M & = \sum_{G \in \calL(M)} x^{-\rank(G)}\left(\sum_{F \geq G} x^{2\rank(F) - 2\rank(E)}\mu[F, E]\cdot \hat{Q}_{M_G^F}(x^{-2})\right)\cdot \zeta^G \\
          & = \sum_{G \in \calL(M)} x^{-\rank(G)}\hat{Y}_{M/G}(x^{-2})\cdot \zeta^G                                                                        \\
\end{align*}

\section*{Computing $\Delta(\mu_M)$}

We defined $\mu_M = \sum_{F \in \calL(M)} x^{\rank(F) - 2\rank(E)}\mu[F,E]\cdot F$. Then,
\begin{align*}
    \Delta(\mu_M) & = \sum_{F \ni e \wedge F\setminus e \in \calL(M)}x^{\rank(F) - 2\rank(E)-1}\mu[F,E]\cdot (F\setminus e)
    + \sum_{\text{other terms}} x^{\rank(F) - 2\rank(E)}\mu[F,E]\cdot (F\setminus e)                                        \\
                  & = \sum_{F \in S}x^{\rank(F) - 2\rank(E)}\mu[F \cup e,E]\cdot F + \mu_{M\setminus e}
\end{align*}

The desired expression for $\Delta(\mu_M)$ is $\mu_{M\setminus e}$ which we could not obtain.

\end{document}